论文写作过程中最重要的是{\bf 图、表、公式}等内容的编排\footnote{这是一个注脚，可以输入你对内容的补充说明}。

%=========================================================================================
\BiSection{图}{Figures}

图~\ref{fig_ch2}~是用~Tikz~画的，Visio~画不出这么好看的图。
%
\begin{figure}[!ht]
\centering
\includegraphics[width=\textwidth]{echoes}
\caption{图注是五号字。} \label{fig_ch2}
\end{figure}

%=========================================================================================
\BiSection{表}{Tables} 

表格要求采用三线表，如表~\ref{tab_ch2}~所示。
%
\begin{table}[!h]
	\renewcommand{\arraystretch}{1.2}
	\centering\wuhao
	\caption{表题也是五号字} \label{tab_ch2} \vspace{2mm}
	\begin{tabularx}{\textwidth} { 
   >{\centering\arraybackslash}X 
   >{\centering\arraybackslash}X 
   >{\centering\arraybackslash}X 
   >{\centering\arraybackslash}X }
	\toprule[1.5pt]
		Interference & DOA (deg) & Bandwidth (MHz) & INR (dB) \\
	\midrule[1pt]
		1 & $-30$ & 20 & 60 \\
		2 & 20 & 10 & 50 \\
		3 & 40 & 5 & 40 \\
	\bottomrule[1.5pt]
	\end{tabularx}
\end{table}



%=========================================================================================
\BiSection{公式}{Equations}
\BiSubsection{单个公式}{Equations}

单个公式的编号如式~(\ref{equ_ch2_pdf})~所示，该式是标准正态分布的概率密度函数~\upcite{Manolakis2005}，从公式可知~\LaTeX~排版的公式比~MathType~美观，而且公式编写效率更高。
%
\begin{equation} \label{equ_ch2_pdf}
f_Z(z) = \frac{1}{\pi\sigma^2} \exp\left(-\frac{|z-\mu_Z|^2}{\sigma^2}\right)
\end{equation}

%-----------------------------------------------------------------------------------------
\BiSubsection{多个公式}{Subequations}

多个公式作为一个整体可以进行二级编号，如~(\ref{equ_ch2_fourier})~所示，该式是连续时间~Fourier~变换的正反变换公式~\upcite{Vetterli2014}。
%
\begin{subequations} \label{equ_ch2_fourier}
\begin{align}
X(f) &= \int_{-\infty}^{\infty}x(t)e^{-j2\pi f t}\dif t \\
x(t) &= \int_{-\infty}^{\infty}X(f)e^{j2\pi f t}\dif f
\end{align}
\end{subequations}

% !Mode:: "TeX:UTF-8" 

\BiChapter{论文结构}{Structure}

\BiSection{摘要}{Abstract}

摘要是论文的高度概括，是全文的缩影，是长篇论文不可缺少的组成部分。要求用中、英文分别书写，一篇摘要不少于 400 字。英文摘要与中文摘要的内容和格式必须一致。

\BiSection{主要符号表}{Symbol Table}

如果论文中使用了大量的物理量符号、标志、缩略词、专门计量单位、自定义名词和术语等，应将全文中常用的这些符号及意义列出\footnote{\url{https://es.overleaf.com/learn/latex/List_of_Greek_letters_and_math_symbols}}。 \textbf{如果上述符号和缩略词使用次数不多，可以不设专门的主要符号表，但在论文中出现时须加以说明。}论文中主要符号全部采用法定单位，严格执行国家标准（GB3100～3102—93）有关“量和单位”
的规定。单位名称采用国际通用符号或中文名称，但全文应统一，不得两种混用。 缩略词应列出中英文全称。

\BiSection{正文}{Main body}

\BiSubsection{绪论}{Introduction}

绪论相当于论文的开头，它是三段式论文的第一段（后二段是本论和结论）。绪论
与摘要写法不完全相同，摘要要写得高度概括、简略，绪论可以稍加具体一些，文字以 1000 字左 右为宜。绪论一般应包括以下几个内容： 

(1)为什么要写这篇论文，要解决什么问题，主要观点是什么。

(2)对本论文研究主题范围内已有文献的评述（包括与课题相关的历史的回顾，资料来源、性质
及运用情况等）。

(3)说明本论文所要解决的问题，所采用的研究手段、方式、方法。明确研究工作的界限和规模。
概括论文的主要工作内容。

\BiSubsection{课题的研究方法与手段}{Methodology}
课题研究的方法与手段，分别以下面几种方法说明： 用实验方法研究课题，应具体说明实验用的装置、仪器、原材料的性能等是否标准，并应对所
有装置、仪器、原材料做出检验和标定。对实验的过程和操作方法，力求叙述得简明扼要，对实验 结果的记录、分析，对人所共知的或细节性的内容不必过分详述。 用理论推导的手段和方法达到研究目的的，这方面内容要精心组织，做到概念准确，判断推理
符合客观事物的发展规律，要做到言之有序，言之有理，以论点为中心，组织成完整而严谨的内容 整体。
用调查研究的方法达到研究目的的，调查目标、对象、范围、时间、地点、调查的过程和方法
等，这些内容与研究的最终结果有关系，但不是结果本身，所以一定要简述。但对调查所提的样本、 数据、新的发现等则应详细说明，这是结论产生的依据。

\BiSubsection{结论与展望}{Conclusion}

在写作时，应对研究成果精心筛选，把那些必要而充分的数据、现象、样品、认识等选出来，
写进去，作为分析的依据，应尽量避免事无巨细，把所得结果和盘托出。在对结果做定性和定量分 析时，应说明数据的处理方法以及误差分析，说明现象出现的条件及其可证性，交代理论推导中认 识的由来和发展，以便别人以此为根据进行核实验证。对结果进行分析后所得到的结论和推论，也 应说明其适用的条件和范围。恰当运用表和图作结果与分析，是科技论文通用的一种表达方式。

结论与展望：结论包括对整个研究工作进行归纳和综合而得出的总结；所得结果与已有结果的
比较；联系实际结果，指出它的学术意义或应用价值和在实际中推广应用的可能性；在本课题研究 中尚存在的问题，对进一步开展研究的见解与建议。结论集中反映作者的研究成果，表达作者对所 研究课题的见解和主张，是全文的思想精髓，是全文的思想体现，一般应写得概括、篇幅较短。

\BiSubsection{致谢}{Thanks}

对于毕业设计（论文）的指导教师，对毕业设计（论文）提过有益的建议或给予 过帮助的同学、同事与集体，都应在论文的结尾部分书面致谢，言辞应恳切、实事求是。


\BiSection{参考文献与注释}{Reference}

参考文献是为撰写论文而引用的有关文献的信息资源。
参考文献列示的内容务必实事求是。论文中引用过的文献必须著录，未引用的文献不得虚列。
遵循学术道德规范，杜绝抄袭、剽窃等学术不端行为。
参考文献须是作者亲自考察过的对学位论文有参考价值的文献。
参考文献应有权威性，应注意所引文献的时效性。 参考文献的数量：一般不少于 10 篇，其中，期刊文献不少于 8 篇，国外文献不少于 2 篇，均以近 5 年的文献为主。 注释是正文需要的解释性、说明性、补充性的材料、意见和观点等。


参考文献格式应符合国家标准 GB/T-7714-2005《文后参考文献著录规则》。中国国家标准化管理委员会于2015年5月15日发布了新的标准 GB/T 7714-2015《信息与文献参考文献著录规则》。因为二者的差别非常小，所以采用了新的标准。标准的 BiBTeX 格式网上资源非常多，本文使用了李泽平开发的版本~\upcite{Lee2016}。本模板中使用 bib\LaTeX  提供参考文献支持。请使用\texttt{\textbackslash upcite} 命令进行引用。

如，本文参考了如下文献\upcite{zhou2020egoplanner} \upcite{thrun2002probabilistic}。



\BiSection{装订要求}{Bindings}

论文的查重部分和不查重部分的内容如表\ref{tab_check}所示。

\begin{table}[!h]
	\renewcommand{\arraystretch}{1.2}
	\centering\wuhao
	\caption{论文查重部分与不查重部分内容对照表} \label{tab_check} \vspace{2mm}
	\begin{tabularx}{\textwidth} { >{\centering\arraybackslash}X >{\centering\arraybackslash}X }
	\toprule[1.5pt]
		查重部分 & 不查重部分 \\
	\midrule[1pt]
		封面        &   致谢                \\
		中文摘要    &   外文原文（附录）    \\
		英文摘要    &   外文译文（附录）    \\
		目录        &   有关算法（附录）    \\
		正文        &   有关图纸（附录）    \\
		参考文献    &   计算机源程序（附录）\\    
	\bottomrule[1.5pt]
	\end{tabularx}
\end{table}

\noindent \textbf{毕业设计(论文)装订次序要求} \\
第一 ~ 封面 \\
第二 ~ 任务书（双面打印） \\
第三 ~ 考核评议书（背面是答辩结果） \\
第四 ~ 中文摘要 \\
第五 ~ 英文摘要 \\
第六 ~ 目录 \\
第七 ~ 正文（含绪论和结论） \\
第八 ~ 致谢 \\
第九 ~ 参考文献 \\
第十 ~ 附录（含外文原文及其译文、有关图纸、计算机源程序等）

% !Mode:: "TeX:UTF-8" 

\BiChapter{格式要求}{Formats}

\BiSection{封面}{Cover}

采用西安交通大学毕业设计（论文）统一封面（模板见毕业设计管理系统——资料下载）， 用 A3 纸（封面封底连在一起，从左侧包住论文）。封面上所填内容居中排列三字号加粗，论 文题目不能超过 35 个汉字。封面可去教材科购买。


\BiSection{任务书、考核评议书}{}

双面打印。《考核评议书》、《评审意见书》和《答辩结果》必须分别由 指导教师、评阅人和答辩组据实填写。

\BiSection{中、英文摘要页}{Abstracts}

摘要页由摘要正文、关键词等组成。

摘要／ABSTRACT 按一级标题编排中文“摘要”二字，二字间距为两个字符。英文为“ABSTRACT”。摘要正文，中文每段开头左起空两字符起排，段与段之间不空行；英文每段开头左对齐顶格编
排，段与段之间空一行。小四号字。 

关键词/KEY WORDS” 正文内容下，空一行，左对齐顶格编排“关键词/KEY WORDS”（小四号，加粗），后接冒号， 其后为具体关键词（小四号），每一关键词之间用分号分开，最后一个关键词后无标点符号。英文 每组KEY WORDS 的第一个字母为大写，其余为小写。

资助申明：资助申明编排在摘要第一页的页脚处。

\BiSection{中、英文目录页}{Contents}
中文目录页应放在奇数页上起排。 “目录”二字按一级标题编排，两字间距两个字符。 目录正文，包括编号、标题及其开始页码。一般只列到三级标题。目录中标题的编号应与正文
中标题的编号一致； 第一级标题左对齐顶格编排；与上一级标题相比，下一级标题左缩进一个字符起排； 标题与页码之间用“……”连接。页码右对齐顶格编排； 建议采用文本编辑软件的“目录自动生成功能”生成目录。 如果有英文目录，英文目录的内容、格式均须与中文目录一致。

\BiSection{表格}{Tables}
图、表、公式等的序号用阿拉伯数字分章连续编号，如图 \ref{fig_ch2}、表 \ref{tab_ch2}、表 \ref{tab_check} 等，但
不出现“公式"两字，将编号置入小括号中，如（\ref{equ_ch2_fourier}）等。图、表和公式等与正文之间间隔 0.5 行。 图应有图题，表应有表题，并分别置于图号和表号之后。图号和图题置于图下方的居中位置，
表号和表题应置于表上方的居中位置。引用图或表应在图题或表题右上角标出文献来源。 若图或表中有附注，采用英文小写字母顺序编号，附注写在图或表的下方。 物理量及量纲均按国际标准（SI）及国家规定的法定符号和法定计量单位标注，禁止使用已废
弃的符号和计量单位。物理量的符号由斜体字母标注，单位的符号使用正体字母标注，量与单位间 用斜线隔开。例如： $I/A$，$\rho/kg \cdot m^{-3}$，$F/N$，$v/m \cdot s^{-1}$ 等等。

表格应紧跟文述编排。表格中一般是内容和测试项目由左至右横读，数据依序竖读，应有自明 性。表的各栏均应标明“量或测试项目、符号、单位”。只有在无必要标注的情况下方可省略。表内 同一栏的数字必须上下对齐。表内不宜用“同上”、“同左”和类似词，一律填入具体数字或文字。表内“空 白”代表未测或无此项，“…”代表未发现，“0”代表实测结果确为零。如数据已绘成曲线图，可不再列 表。

\begin{enumerate}[wide, label=\arabic*),  labelindent=21pt]
    \item 表格转页接排时，在随后的各页上应重复表的编号。编号后跟（续），如表 l（续），续表
均应重复表头和关于单位的陈述。
    \item 一律使用三线表，与文字齐宽，上下边线，线粗 1.5 磅，表内线，线粗 1 磅。在三线表中可以加辅助线，以适应较复杂表格的需要。
    \item 使用他人表格须注明出处。
    \item 表中用字一般为五号字。如排列过密，用五号字有困难时，可小于五号字，但不小于七号
\end{enumerate}


\BiSection{图片}{Figures}

\begin{enumerate}[wide, label=\arabic*),  labelindent=21pt]
    \item 一幅图如有若干幅分图，均应编分图号，用(a),(b),(c),…按顺序编排；
    \item 插图须紧跟文述。在正文中，一般应先见图号及图的内容后见图，特殊情况须延后的插图
不应跨节； 
    \item 提供照片应大小适宜，主题明确，层次清楚，利于复制，金相照片一定要有比例尺；
    \item 图应具有"自明性”，即只看图、图题和图例，不阅读正文，就可理解图意。
    \item 图中的标目是说明坐标轴物理意义的项目，由物理量的符号或名称和相应的单位组成。
    \item 图中用字一般为五号字，如排列过密，用五号字有困难时，可小于五号字，但不得小于七
号字。 
    \item 图的大小一般为宽 6.67cm $\times$ 高 5.00cm。特殊情况下，也可宽 9.00cm $\times$ 高 6.75cm，或宽 13.5cm
$\times$ 高 9.00cm。同类图片的大小应一致。图片的编排应美观、整齐。
\end{enumerate}


\BiSection{其他}{Others}

\BiSubsection{页码}{}

论文页码的第一页\textbf{从正文开始}用\textbf{阿拉伯数字}标注，直至全文结束。
\textbf{正文前的内容}（除 封面）用\textbf{罗马数字}单独标注页码。
页码位于页面底端，对齐方式为 “外侧”，页码格式为最简单的数字，不带任何其它的符号或信息。
页码不能出现缺页和重复页。附录（含外文原文及其译文、有关图纸、计算机源程序等）必须与论文装订在一起，附录的页码必须接着参考文献的页码连续编写。